problem:  
Let \( (M^n, g) \) be a complete, noncompact Riemannian manifold satisfying the following:  
1. The sectional curvature \( \mathrm{sec}_M \ge 0 \) outside a compact subset \( K \subset M \);  
2. \( M \) has **Euclidean volume growth**, i.e.  
   \[
   \lim_{r \to \infty} \frac{\mathrm{Vol}(B_r(p))}{r^n} = v_0 > 0;
   \]
3. Every tangent cone at infinity of \(M\) (in the pointed Gromov–Hausdorff sense) is isometric to a fixed product \( C(Z) \times \mathbb{R}^k \) for some compact Alexandrov space \( Z \) with \( \mathrm{curv}(Z) \ge 1 \).

**Question:**  
Under these hypotheses, does the asymptotic structure \( C(Z)\times\mathbb{R}^k \) uniquely determine the **topology of the end** of \( M \)?  
That is, if \( (M_1, g_1) \) and \( (M_2, g_2) \) are two such manifolds whose tangent cones at infinity are both isometric to \( C(Z)\times\mathbb{R}^k \), must their complements of compact subsets be **homeomorphic (or diffeomorphic) to each other**?  

Equivalently, does the uniqueness of the tangent cone at infinity, together with Euclidean growth and eventual nonnegative curvature, enforce **end rigidity**: that manifolds sharing the same asymptotic cone have the same large-scale topological type?

Why is it a "good" problem:  
This problem sits squarely at the joint between **asymptotic metric geometry** and **topological rigidity theory**. It asks whether the global topological structure near infinity is already encoded by the metric cone one sees from infinitely far away—a strikingly simple but unproven principle.  

The key insight is that, for nonnegatively curved manifolds, tangent cones at infinity describe the “macroscopic profile” of the manifold, while the Soul Theorem describes its compact core. If the end topology is rigidly dictated by the cone’s geometry, one would arrive at an **asymptotic version of the Soul Theorem**, converting metric asymptotics into topological uniqueness.  

The problem is clear yet deep: it brings together tools from **Cheeger–Colding theory** (metric limits and cone stability), **Alexandrov geometry** (curvature bounds in nonsmooth spaces), and **Riemannian topology** (structure of ends). A resolution would create a new bridge between asymptotic geometry and global topology, and establish a new form of rigidity parallel to the classical nonnegative-curvature rigidity phenomena. Its simplicity, precision, and potential for foundational impact make it a “good” research problem.